\subsection{Decisions and Strategies}

Agents decide whether to engage in a speculative attack (i.e., acquire foreign currency for exit) based on their posterior beliefs about the fundamental $\theta$.

Following the global-games literature, we restrict attention to monotone strategies. There exists a cutoff $x^*$ such that agent $i$ attacks if and only if:
\begin{equation}
x_i \geq x^*.
\end{equation}

\subsection{Aggregate Attack}

Given the cutoff strategy, the measure of agents who attack is:
\begin{equation}
R(\theta) = \mathbb{P}(x_i \geq x^* \mid \theta) = 1 - \Phi\left(\frac{x^* - \theta}{\sigma_x}\right),
\end{equation}
where $\Phi(\cdot)$ is the standard normal cumulative distribution function.

\subsection{Hedging Demand}

In addition to speculative demand, agents demand foreign currency for hedging purposes. We model aggregate hedging demand as an increasing function of expected misalignment:
\begin{equation}
H = h \cdot \mathbb{E}[\theta \mid s], \quad h > 0.
\end{equation}

This captures the idea that even in the absence of coordination, agents acquire foreign currency as insurance against depreciation risk.

\subsection{Crisis Condition}

A crisis occurs if total demand for foreign currency exceeds the available policy buffer. We model this as:
\begin{equation}
R(\theta) + H \geq \alpha - \delta \theta,
\end{equation}
where:
\begin{itemize}
\item $\alpha$ represents the strength of policy buffers (e.g., reserves)
\item $\delta$ captures the fact that higher misalignment weakens the regime
\end{itemize}

\subsection{Crisis Probability}

Given the information structure, agents form beliefs about the probability of a crisis:
\begin{equation}
\pi(x_i, s) = \mathbb{P}\left(R(\theta) + H \geq \alpha - \delta \theta \mid x_i, s\right).
\end{equation}

\subsection{Payoffs}

Agents derive utility from consumption, which depends on whether a crisis occurs.

\paragraph{No-crisis consumption}
\begin{equation}
c^N = y + \zeta \log(1 + q_c + q_s) - C_c(q_c) - C_s(q_s),
\end{equation}

\paragraph{Crisis consumption}
\begin{equation}
c^C = c^N - \lambda \theta (1 + \xi \omega_s),
\end{equation}

where $\omega_s = \frac{q_s}{q_c + q_s}$ is the share of stablecoins in total foreign currency holdings.

\paragraph{Expected utility}
\begin{equation}
U = (1 - \pi) \log(c^N) + \pi \log(c^C).
\end{equation}

\subsection{Equilibrium Definition}

An equilibrium consists of:
\begin{itemize}
\item A cutoff $x^*$
\item Asset allocations $(q_c, q_s)$
\item Beliefs consistent with Bayesian updating
\item A crisis probability consistent with aggregate behavior
\end{itemize}

such that agents optimize given beliefs and the information structure, and all objects are mutually consistent.

\subsection{Key Mechanism}

Stablecoins affect equilibrium through two channels:
\begin{enumerate}
\item \textbf{Cost channel:} lower transaction costs increase foreign currency demand
\item \textbf{Information channel:} improved public signal strengthens coordination
\end{enumerate}

The second channel amplifies strategic complementarities and increases the likelihood of coordinated attacks, generating fragility in equilibrium.
